% ============================================================
% CHAPTER 11: CHALLENGES & MITIGATIONS
% MCDA 5585 / 5586 Major Project Report
% Bhavik Kantilal Bhagat | A00494758 | Saint Mary's University
% ============================================================

\chapter{Challenges \& Mitigations}
\label{chap:challenges}

Six engineering challenges arose during my internship that were significant
enough to alter the delivery plan or require a design revision. Each is
documented with the problem encountered, the mitigation I applied, and the
tradeoffs accepted.

% ────────────────────────────────────────────────────────────
\section{Challenges and Mitigations}
\label{sec:challenges}
% ────────────────────────────────────────────────────────────

Six engineering challenges arose during my internship that were significant
enough to alter the delivery plan or require a design revision. They are
documented here both for completeness and as a record of the problem-solving
approach I applied at each decision point.

\subsection{Challenge 1: CloudWatch Metric Insights Percentile Limitation}

The \texttt{GROUP BY} mode of CloudWatch Metric Insights SQL supports only
\texttt{AVG}, \texttt{SUM}, \texttt{MIN}, and \texttt{MAX} aggregate functions;
\texttt{p50} and \texttt{p99} percentile statistics are not available in this
mode. This became apparent when I attempted to build a Lambda duration latency
distribution panel that would show P50 and P99 across all 119 functions in a
single query. The mitigation I adopted was a wildcard dimension query in standard
CloudWatch metrics mode (\texttt{FunctionName: ["*"]}), which returns individual
P50/P99 time series for each recently active function. The accepted tradeoff is
that functions with no recent invocations are absent from the panel --- a gap
I documented as a known limitation and flagged as a future-work item for
custom-dimension metric emission.

\subsection{Challenge 2: Grafana 10.4 $\to$ 12.4 Mid-Development Upgrade}

Two days after my initial CloudFormation stack deployment on May~11, the
Amazon Managed Grafana workspace was upgraded from version~10.4 to version~12.4.
This was a major version migration that introduced breaking changes in the
Metric Insights SQL query editor and in template variable interpolation syntax.
Approximately one full working day was consumed auditing all 110+~panels,
identifying the affected queries, and migrating them to the stable Grafana~12.4
CloudWatch query model. The experience produced NFR6: all future panel queries
I author would use only the standard CloudWatch plugin query interface, avoiding any
AMG-version-specific syntax that would break on the next major upgrade.

\subsection{Challenge 3: Secrets Manager 64\,KB Limit}

The CloudFormation IaC provisioner I built stored its YAML alarm templates as a
single JSON-encoded secret in AWS Secrets Manager. As the monitoring scope
expanded to cover all seven service categories, the generated YAML grew to
\textbf{14,720 lines at version v728}, exceeding the hard 64\,KB limit on
Secrets Manager secret values. The Lambda function began failing with
\texttt{SecretSizeExceeded} on the \texttt{PutSecretValue} call. The
mitigation was a migration of all template files to S3 with object versioning
enabled, providing unlimited payload size, native rollback capability, and
negligible additional cost. I updated the Lambda execution role to include
\texttt{s3:GetObject} on the template bucket, a minimal IAM surface expansion
that was accepted in exchange for the elimination of all payload size
constraints going forward.

\subsection{Challenge 4: Shared BudgetCode Tag Ambiguity}

The \texttt{CITCOWORKS} and \texttt{MERIDIAN} environments both carry
\texttt{BudgetCode=``Shared Cloud Artifacts''} --- a legacy of both platforms
being onboarded before the IA team standardized the uniqueness requirement for
\texttt{BudgetCode} values. A single-stage tag filter could not distinguish
resources from the two environments, causing inventory contamination in early
prototype runs of the inventory API. Retroactively retagging all production
resources across two platforms was ruled out as too disruptive. My mitigation
was the \textbf{two-stage \texttt{AppName} filter} documented in
Section~\ref{sec:inventory_api}: any discoverer operating in the
\texttt{``Shared Cloud Artifacts''} BudgetCode must apply a secondary
\texttt{AppName} check as a mandatory discriminator. This created a tagging
governance obligation --- new resources added without the correct
\texttt{AppName} tag are invisible to the inventory API --- mitigated by a tag
compliance CloudWatch alarm and a documented platform-wide constraint in the
repository \texttt{README}.

\subsection{Challenge 5: IAM Permission Boundary Blocking Subscription Filters}

The Top-K centralized log aggregation design evaluated under IISS-469 required
\texttt{logs:PutSubscriptionFilter} to route log events from all IA Lambda
functions into a central processing pipeline. Systematic analysis determined
that the approach would require \textbf{106 subscription filters} at the
current fleet size. The attempt to provision these filters programmatically
surfaced an IAM permission boundary constraint: the IA team's scoped permission
set did not include \texttt{logs:PutSubscriptionFilter}, and elevating it
required an approval process that would have consumed disproportionate
calendar time relative to the benefit. The design was redirected toward the
CloudWatch Metric Insights SQL approach, which I used to achieve comparable per-resource
drill-down without requiring subscription filter infrastructure.

\subsection{Challenge 6: Module-Import Credential Caching (IISS-706 Hotfix)}

The CAIS Deadline hotfix on June~8 exposed a fragility in the credential
retrieval pattern used by the affected bot: credentials were fetched once at
module import time and cached for the life of the process. When AWS Secrets
Manager underwent a scheduled rotation during an active bot execution, the
rotation returned an empty credential set during the brief transition window.
Because the module-level cache had already been populated from the previous
(valid) secret version, the bot continued executing with what it believed were
valid credentials --- and encountered authentication failures only when
attempting the actual UCM login. The fix required me to restructure the credential
retrieval to occur immediately before each use rather than at import time,
ensuring that a fresh \texttt{GetSecretValue} call is made on each execution
cycle. I documented this pattern as a platform-level recommendation for all
RPA workflows that interact with Secrets Manager.
